\documentclass[instructions.tex]{subfiles}
\begin{document}

\chapter{Le péché afflige Dieu et le déshonore.}


\primo Il faut que l'injure que le péché fait à Dieu lui soit bien sensible et bien insupportable, puisque l'Écriture nous apprend que Dieu, voyant la malice des hommes, se repentit de les avoir formés, et fut pénétré de douleur jusqu'au fond du cœur\footnote{Tactus dolore cordis intrinsecùs. Gen.6.}. Ce n'est pas que Dieu puisse ressentir de la douleur, puisqu'il est immuable; mais ces paroles de l'Écriture nous marquent combien le péché est opposé à Dieu; si opposé, que si Dieu pouvait souffrir en lui-même, le péché, autant qu'il est en lui, serait capable de le faire souffrir et de l'anéantir, et qu'un seul péché mortel lui ferait ressentir plus de tristesse que les bonnes œuvres de tous les justes ne pourraient lui causer de joie.

Ce que le péché a fait souffrir à un Dieu dans l'humanité du Verbe incarné, en est une preuve. Qu'est-ce qui a accablé cet Homme-Dieu d'une tristesse mortelle, jusqu'à lui faire suer le sang? n'est-ce pas la vue des péchés du monde? Oui, dit un Père, la vue du péché a plus causé de douleur au Fils de Dieu que les plaies de son corps\footnote{Plus gravant vulnera peccati tui, quàm vulnera corporis sui.}. Que votre cœur est dur, si vous aimez le péché, si vous prenez plaisir d'affliger ainsi votre Dieu! Comprenez du moins, par l'affliction qu'il lui cause, quelle douleur vous devez vous-même en ressentir.


\secundo Le dernier trait de la malignité du péché, c'est qu'il déshonore Dieu et lui ravit sa gloire\footnote{Vos inhonorâstis me. Joan.8.}. Gloire due à Dieu seul, et qui lui est si propre, qu'il ne peut y renoncer; gloire qui est le seul bien qu'il puisse tirer de sa créature. Dieu ne peut être glorifié en lui-même que par lui seul; mais il était convenable que ses perfections et ses grandeurs fussent glorifiées au dehors. C'est pour cette fin qu'il a formé l'univers, qu'il le remplit de la majesté de sa gloire, et qu'il a produit des créatures intelligentes, capables de le connaître, pour leur manifester et leur faire glorifier ses grandeurs.

Voilà ce que ne connaît pas l'homme pécheur, et ce qu'il refuse à Dieu. Il oublie ce qu'il doit à sa grandeur, en se comportant comme si Dieu était indigne de ses hommages.

Dieu n'est ni moins puissant ni moins heureux, il est vrai, mais il est moins honoré par sa créature; honneur qu'elle lui doit, et pour lequel il a formé cette créature. Refuser à Dieu cet honneur, c'est donc le priver d'une chose qui lui est due; ce qui est toujours un mal: et plus Dieu est grand, plus aussi ce mal est grand, et plus l'homme qui lui fait cette injure est criminel.

Un roi mérite les hommages de ses sujets; et, quoique la puissance d'un roi ne soit pas affaiblie parce qu'un vil esclave lui manque de respect, ce vil esclave n'en commet pas moins un crime envers son prince. Plus il est vil, plus aussi il est criminel et punissable de manquer à ce qu'il doit à son souverain.

Dieu est infiniment adorable, c'est l'Être suprême: toutes les créatures réunies sont si peu de chose et si viles à ses yeux, qu'elles sont comme rien, dit l'Écriture\footnote{Omnes gentes quasi non sint. Is.40.}. Or, que cette majesté souveraine, devant qui tous les monarques ne sont que poussière, soit méprisée par une vile créature qui devrait l'adorer jusqu'à l'anéantissement, s'il était possible, n'est-ce pas là le comble de l'outrage? outrage si grand, que tous les hommes, quand ils s'offriraient en sacrifice, ne pourraient jamais le réparer par eux-mêmes. Non, tout ce que les justes ont jamais fait de plus héroïque, ne donnera jamais tant de gloire à Dieu, qu'il en perd par le péché mortel d'une seule créature; parce que les hommages des justes sont finis et limités, au lieu que l'outrage que Dieu reçoit par le péché mortel est comme infini, par cette raison que l'énormité d'un outrage se mesure par la bassesse de celui qui offense et par la grandeur de celui qui est offensé. L'injure que le péché mortel fait à Dieu est donc comme infinie, puisqu'elle offense un être infini. Cette injure ne peut donc être réparée par les hommages de tous les hommes.

C'est pour cette fin, pour réparer la gloire de Dieu son Père, que le Fils de Dieu s'est anéanti jusqu'à se faire victime de nos péchés, parce qu'il n'y a qu'un Homme-Dieu qui puisse honorer Dieu autant qu'il le mérite, et réparer les outrages faits à sa majesté. Saint Augustin avait donc bien raison de dire qu'il vaudrait mieux que le ciel et la terre, et tout ce qui n'est pas Dieu pérît, que d'offenser Dieu, même légèrement\footnote{cf. S. Aug.}. La destruction de l'univers ne serait que le mal de la créature, au lieu que le péché est le mal du Créateur. Or, le moindre mal du Créateur est au-dessus de tous les maux des créatures. Oh! que le péché est donc bien à craindre! et que celui-là est malheureux, qui ne le comprend pas!


\section{Résolutions}

Il faut que je travaille dès aujourd'hui à bannir le péché de toute ma conduite.

\primo Je m'attacherai donc à considérer l'horreur que Dieu en a, et combien ce monstre lui ravit de gloire.

\secundo Ah! mon Dieu, vous m'avez créé pour vous, et jusqu'ici j'ai presque toujours vécu contre vous; convertissez-moi, pardonnez-moi.

\end{document}
